\section{The mechanism}
\label{sec:mechanism}

The model keeps two questions apart and answers them separately, because
conflating them is what makes long-run fertility arguments irresolvable:

\begin{description}[leftmargin=1.4em,style=nextline]
  \item[Who is becoming more common?] The composition of the population by
  fertility disposition. This is what selection moves.
  \item[How many children does a person of a given disposition have?] The
  fertility environment. This is what social and economic change moves.
\end{description}

Observed fertility is the product of the two. It is an output of the model, not
an input to it, and that is the structural difference from a conventional
projection, in which observed fertility is the modelled object itself.

\subsection{Types, propensities and transmission}

Each country's population is divided along an extra axis into $K$ fertility
\emph{types}. A type is not an ethnic, religious or political label; it is an
effective fertility propensity $\propensity_k$, a multiplier on the country's own
base-year age-specific fertility schedule. A woman of propensity $1.4$ has forty
per cent more children than the national average did in the base year, at every
age, and that multiplier does not change when conditions do. The propensities are
normalised so that the base-year composition reproduces each country's published
fertility exactly, which anchors the model to the present instead of fitting it
to the future.

Writing $\ell$ for country, $t$ for year, $a$ for single year of age, $k$ for
type, $f_t[\ell,a]$ for the baseline age-specific fertility schedule and
$\envmult_t[\ell]$ for the environment multiplier, births are counted by the
mother's type and then reassigned to the child's:
\begin{align}
  B_t[\ell,k] &= \sum_a \bar{P}_t[\ell,k,\mathrm{F},a]\;
                 f_t[\ell,a]\;\envmult_t[\ell]\;\propensity[\ell,k],
                 \label{eq:births-by-mother}\\[0.2em]
  \tilde{B}_t[\ell,j] &= \sum_k B_t[\ell,k]\; M_t[\ell,k,j],
                 \label{eq:births-by-child}
\end{align}
where $\bar{P}_t$ is mid-year exposure of women and $M_t$ is the transition
matrix from mother's type to child's type. Everything the mechanism claims is in
those two lines. Survival, ageing, the open-ended age group and migration are
untouched and carried identically for every type.

Selection is what happens when the rows of $M$ differ and $\propensity$ is not
constant. Higher-propensity women appear more often in
\Cref{eq:births-by-mother} because they have more children; those children land
in their mother's type more often than chance through
\Cref{eq:births-by-child}; so the composition drifts upward, the
population-average propensity rises, and observed fertility departs from the
environment path without anyone imposing a floor on it. The quantity we report
throughout is that departure,
\begin{equation}
  S_t[\ell] \;=\; \frac{\text{observed fertility}}{\text{environment alone}}
            \;=\; \bar{\propensity}_t[\ell],
  \label{eq:selection-effect}
\end{equation}
the propensity averaged over women of childbearing age. It equals one when
selection is switched off, and its distance above one is the size of the
mechanism's claim.

\subsection{The response, and what it does and does not rest on}
\label{sec:response}

The mainstream transition rule is deliberately the simplest one that can carry
the evidence. A child takes its mother's type with probability $\rho$ and is
otherwise drawn from the surrounding population of childbearing adults. Under
that rule the one-generation response has a closed form. Mothers are sampled in
proportion to their own fertility, so the birth-weighted distribution of types is
$q_k \propto w_k \propensity_k$ where $w_k$ is the share of childbearing women of
type $k$. The mean propensity of the newborn cohort is then
\begin{equation}
  \rho \sum_k q_k \propensity_k + (1-\rho) \sum_k w_k \propensity_k
  \;=\; \bar{\propensity} + \rho\,\frac{\sigma^2}{\bar{\propensity}},
  \label{eq:response}
\end{equation}
so the relative gain per generation is $\rho\,\mathrm{CV}^2$: the transmission
coefficient times the squared coefficient of variation of family size. This is
the secondary theorem of natural selection in its covariance form
\citep{robertson1966,price1970} --- the response equals the covariance of the
trait with fitness divided by mean fitness --- specialised to the case where the
trait and the fitness are the same quantity, so that the covariance is a
variance.

Three things about that equation need to be said plainly, because each is an
assumption that could be wrong.

\textbf{The transmission probability and the reported correlation coincide only
because the variances are equal.} Under this rule a child is either identical to
its mother in type or an independent draw from the same distribution, so the
implied parent--offspring correlation in propensity is exactly $\rho$, and
$\mathrm{Cov} = \rho\sigma^2$ rather than the general $\rho\sigma_X\sigma_Y$.
That is a property of the rule, and it is why $\rho$ can be read off a reported
correlation. It also means the model assumes the dispersion of family size does
not change between generations, which is a real restriction and not a harmless
normalisation.

\textbf{The covariance is constructed, not observed.} The dispersion comes from
one set of countries and cohorts and the correlation from another
(\Cref{sec:calibration}). Multiplying them produces the covariance a single
intergenerational sample would have if both applied to it; it does not reproduce
a covariance measured in one. The right estimate would come from paired
parent--child completed-fertility data with both moments taken from the same
sample, and we did not find that at the country coverage this model needs. We
regard this as the weakest link in the calibration and have not disguised it by
calling the product a measurement.

\textbf{Matching one generation does not validate four.} Many transition rules
reproduce the same first-generation response and compound differently over
\ProjectionEndYear{}. What is checked is the first generation, which the
implementation reproduces to nine decimal places; what is assumed is that the
rule keeps holding. \Cref{sec:limitations} states which alternative rules would
change the answer and by how much.

At the central calibration ($\rho = \Persistence{}$, $\mathrm{CV} =
\FamilySizeCV{}$) the one-generation response is \GenerationResponse{}.
Compounded over the roughly four generations between \PhaseFiveBaseYear{} and
\ProjectionEndYear{}, that closed form implies a selection multiplier near
$1.22$, against the $\RaceSelectionEffect{}$ the full projection produces. The
closed form runs slightly hot because it ignores the age structure that delays
each generation's effect and the migration that mixes compositions between
countries. That the two agree within a few per cent is a check that the
simulation is doing what the theory says.

\subsection{Named groups are transitions, not labels}
\label{sec:named-groups}

Populations with both very high fertility and strong boundaries are the version
of this mechanism that reaches public discussion, and they are modelled here as
an additional type with two properties rather than one. A named group has its own
fertility, and its own \emph{retention}: the share of children raised in it who
remain in it as adults. Those who leave are not deleted into an average; they
land in the mainstream, weighted toward its high-propensity type, on the
reasoning that leaving a community is not the same as abandoning everything it
taught you. How much they carry is a scenario knob (\Cref{sec:calibration}), set
at \DefectorWeight{} and reported separately for that reason.

Retention is as load-bearing as fertility. A group at one per cent of a
population with six children per woman and ninety per cent retention, against a
mainstream at $1.4$, reaches \ToyFinalShare{} of the population in
\ToyGenerations{} generations, while the same group with weak retention goes
nowhere. This is the condition \citet{arenberg2022} state, and it is worth being
careful about its units: theirs is a single-sex condition, that the expected
number of same-type daughters surviving to reproduce exceeds one, not that
retention times the total fertility rate exceeds one. The cohort engine applies
survival and the sex ratio at birth to every cohort, so the condition it enforces
is the single-sex one; the loose phrasing is only ever a shorthand.

Two limits are stated rather than fixed. No conversion \emph{into} a named group
is modelled, which is small for these two groups but not zero. And every type
starts with its host country's age structure, because nobody publishes one by
type; since a population with six children per woman is far younger than its
host, this understates group growth, and every group share reported below is a
lower bound.
